\nirsection{ЗАКЛЮЧЕНИЕ}

В ходе учебной практики «Научно-исследовательская работа» первого семестра выполнен теоретический анализ подходов к управлению структурированными данными, естественно-языковому формированию запросов и адаптации пользовательского интерфейса. Изучены научные публикации, посвящённые семантическому сопоставлению запросов и схем баз данных, использованию обучаемых моделей, организации интеллектуальных информационных систем и оценке пользовательского взаимодействия.

По результатам обзора установлено, что задача создания доступного средства взаимодействия с данными не исчерпывается преобразованием фразы в корректный SQL-синтаксис. Необходимо учитывать семантику предметной области, неоднозначность запросов, различия возможностей СУБД и последствия исполняемых операций. В работах по адаптивным интерфейсам выделяются подходы к учёту характеристик и предпочтений пользователей, однако обоснованность конкретного механизма адаптации требует экспериментальной проверки.

Сформулирована научная проблема согласования трёх характеристик: доступности пользовательского взаимодействия, корректности интерпретации запроса и контролируемости операций над данными. Для последующего исследования предложена концептуальная схема SmartSQL с раздельными компонентами представления метаданных, интерпретации намерения, проверки, исполнения и адаптивного представления результатов. Определены четыре проверяемые гипотезы и предварительные критерии их оценки.

Поставленные задачи теоретического этапа выполнены в объёме анализа литературы и постановки исследования. Практическая реализация, экспериментальная проверка гипотез, количественные результаты и обоснование научной новизны относятся к следующим этапам магистерской работы. Таким образом, сформирована основа для последовательной разработки и проверки интеллектуальной системы управления структурированными данными с адаптивным естественно-языковым интерфейсом.